\chapter{Discussion, Conclusions and Future Work}
\label{ch:conclusions}

\section{The Gap Between Loss and Generation}

The main scientific finding of the direct track is that reducing loss under
teacher forcing is separate from valid autoregressive generation. Version v6.2
sharpens this separation: on the same 32 training samples, the text loss fell by
32.20\%, but the learned behaviour did not transfer to reliable generation. The
explanation consistent with the evidence is exposure bias, not a lack of
capacity. The methodological consequence is that speech-to-speech checkpoint
selection should not rest on offline loss alone.

In the cascade track, two similar separations were seen. First, the 9/9
mechanical panel with the 0.5B responder does not establish the semantic quality
of the 4B responder. Second, passing a same-family automatic test does not
establish perceptual quality. The gap between a character error rate of 7.14\%
and a word error rate of 30.41\% also shows that approximately preserving
characters is not the same as word-level intelligibility.

The low score of the aligned reference in the semantic test is a reminder that the
corpus's next turn is not necessarily an appropriate reply to the user's last
utterance; this is why the LoRA adapter that reduced the development loss did not
beat the baseline.

\section{Why the Direct Track Failed}

The final audit showed that Moshi's English tokenizer needs 6.6 tokens per Persian
word and that 98.5\% of turns exceed the one-token-per-frame budget. Removing this
bottleneck with phoneme romanization raised the best generation result from 3/9 to
6/9 and made silent output audible, but the speech did not become intelligible; so
the tokenizer bottleneck is real but not the only obstacle.
J-Moshi~\cite{ohashi2025jmoshi} also replaced the tokenizer to bring Moshi to
Japanese and then needed about 60,000 hours of dialogue pre-training, more than
five hundred times this project's data. The speech-input Gemma~4 model also showed
that removing recognition, even with the same language model, lowers quality.

\section{Limitations and Considerations}

The latency and interruption evidence is automated and on an H100: the memory of
a 12 GB card was respected but not its speed, and there was no human user, no
physical microphone and loudspeaker path, and no real acoustic echo. The
interruption and data labels are automatic and have no human labels; the two
automatic judges of the responder disagree considerably; elderly speech in the
tests is simulated. Interruption detection relies on the speech recognizer, not
on the acoustic features specified in the project definition. The dataset is
derived from an internal archive and is not redistributable. The results of the
listening study have not yet been received.

\section{Answers to the Research Questions}

\textbf{1.} Low-rank adaptation of Moshika lowered held-out perplexity (text
tokens 2.53 to 1.19, audio 13.98 to 4.71) but did not produce intelligible speech;
the English tokenizer is the main obstacle. \textbf{2.} In the cascade,
speculative end-of-speech detection with audio held until confirmation, streaming
generation with a Persian-script constraint, clause chunking, a playback queue and
two-stage interruption control with heard-only history brought the latency to a
mean of 372 and a 90th percentile of 431 ms. \textbf{3.} Size did not buy quality:
the smallest Persian-specialized recognizers and synthesizer gave the best
trade-off, responders of 9 to 36 billion parameters (four bits where needed)
gave no gain over the best 4B under both judges, and the 7.7-billion-parameter
end-to-end model did not reach intelligible speech. \textbf{4.} The perception of older adults depends on the
results of the listening study. The status of each requirement is in
Table~\ref{tab:req-status}.

\section{Future Work}

Completing the evidence: running the same loopback test on a 12--24 GB card with a
physical audio path; collecting the listening-study responses and a live
interaction session with older adults; human review of a sample of the labels;
and human judging of the responders (Qwen3.5-4B with the elderly prompt has passed
the deployment qualification, and the participants' preference decides on its
replacement). Improving the system: learned separation of listener words from
content words with a clipped onset (a 200 ms pre-roll was tested and rejected),
semantic end-of-turn detection for the pauses of elderly speech (tested in two
rounds; the best gate split fewer utterances but delayed too many finished turns,
so a shorter hold or human-labelled turn ends are needed), and a more accurate recognizer within the same latency
budget. Direct track: a tokenization that fits the frame budget, a remedy for the
gap between teacher forcing and free generation, and the frozen-model approach.
